\chapter{L'interface}

\section{Écran Lecture}
\capture{now-playing-en}{Écran Lecture, informations repliées}

De haut en bas :
\begin{description}
  \item[Barre d'outils] logo, plein écran, réglages en direct, choix de la visualisation (oscilloscope ou réponse
        du filtre), tableau de bord, menu.
  \item[Visualisation] oscilloscope par voix et, en dessous, le tableau de bord de la puce (chapitre~\ref{chap:dashboards}).
  \item[Informations] titre, auteur, puce, format et pilote d'origine. Repliées par défaut pour laisser la place à
        la visualisation ; la flèche en bas à gauche les déplie.
  \item[Position] temps écoulé, barre de progression, durée totale.
  \item[Commandes] informations, aléatoire, précédent, lecture/pause, suivant, répétition, bouton \textbf{$\equiv$+}
        (favoris, ajouter à une liste), menu du morceau (partager, envoyer, sonnerie, propriétés).
\end{description}

\capture{now-playing-info-en}{Écran Lecture, informations dépliées}

\section{Écran Listes}
\capture{playlist-en}{Écran Listes}
En haut, la recherche et les actions sur la liste (statistiques, export, effacer, trier). Sous la recherche, le
sélecteur de liste indique la liste affichée et son nombre de morceaux ; le toucher ouvre le menu des listes
(chapitre~\ref{chap:listes}). Chaque morceau montre la famille de puce, le titre, l'auteur et la durée ; la poignée
à droite permet de le déplacer.

\section{Écran Navigateur}
\capturedouble{browser-sources-en}{browser-folder-en}{Sources}{contenu d'un dossier}
Le chemin parcouru s'affiche en haut ; chaque élément permet de revenir directement à ce niveau.

\section{Mini-lecteur et barre de navigation}
Hors de l'écran Lecture, un mini-lecteur rappelle le morceau en cours et permet de passer au suivant ou de mettre
en pause.

\section{Paysage et tablette}
\capture[0.9]{landscape-now-playing-en}{Écran Lecture en paysage}
\capture[0.9]{tablet-playlist-en}{Sur tablette}

\section{Plein écran}
Le bouton plein écran affiche l'oscilloscope seul, et en portrait le tableau de bord sous lui. En portrait,
l'affichage évite l'encoche de la caméra en haut et les coins arrondis en bas ; en paysage, l'encoche latérale.
Un appui court quitte le plein écran, un appui long ouvre les réglages en direct.
